\begin{helper}[Second-coordinate same-color injection tail]
The first-coordinate same-color fixed-pair injection tail also holds for
marked sets in the second coordinate.  More explicitly, under the same
support, logarithmic, ratio, and base-codegree hypotheses as in
\(\noderef{FiberAndDegreeSameColorLiftedIntersectionFixedPairInjectionTail}\),
the number of injections whose image has more than
\(100(\log n)^3\) points in
\[
\operatorname{Fin}(m)\times
\{t:G_B.Adj(b,t)\land G_B.Adj(c,t)\}
\]
is at most
\[
\exp(-2(\log n)^3)\,
|\operatorname{Inj}(\operatorname{Fin}(n),
\operatorname{Fin}(m)\times\operatorname{Fin}(m))|.
\]
\end{helper}
\begin{proof}
Apply the coordinate-swap bijection \((x,y)\mapsto(y,x)\) to every injected
cell.  This sends injections hitting the second-coordinate marked set above
to injections hitting the first-coordinate marked set
\[
\{t:G_B.Adj(b,t)\land G_B.Adj(c,t)\}\times \operatorname{Fin}(m)
\]
with exactly the same number of marked image points, and it is an involution
on the finite set of injections.  The count is therefore equal to the
first-coordinate count bounded by
\(\noderef{FiberAndDegreeSameColorLiftedIntersectionFixedPairInjectionTail}\).
\end{proof}
